\documentclass[11pt]{article}

\input{0-parameters.tex}

\begin{document}

\begin{abstract}
    Under consistent input--output national accounting (with fixed shares $a_{ji}$), partial substitution of gas-based electricity by renewables in manufacturing does not alter aggregate TFP growth, regardless of whether energy flows are measured in final, useful, or primary form. 
    Using the $\zeta$-framework created in this paper, considering primary, final or useful energy does not change  TFP growth - except when primary energy is analysed within a broad energy sector that includes both electricity and fossil fuels.
Except for this one case: electrification and the switch from fossil to renewables has no impact on productivity
\[
   \hat{A}=0.
\]

In the CCC approach we are unfortunately in the ``exception'' case where a single energy sector sells both MWh of primary fossil fuel and electricity (both from renewables and fossil fuel). \textbf{This is why the best approach is to use useful energy.}
\end{abstract}


\section{Appendix: Switching from Gas to Renewable Electricity (Final, Useful, and Primary)}
\label{appendix_gas_renewables}

Energy accounting often distinguishes three stages of energy flows: \textbf{primary}, \textbf{final}, and \textbf{useful}. By convention:

\begin{itemize} \item \textbf{Primary energy} is energy found in raw fuels or natural sources (coal, crude oil, wind, etc.) \emph{before} any conversion or transformation processes. \item \textbf{Final energy} is the energy commodity \emph{delivered to end-use sectors} after refining, electricity generation, transmission, or other conversion steps. \item \textbf{Useful energy} is the \emph{portion of final energy} that effectively provides energy services (heat, mechanical work, lighting, etc.), after end-use device losses \citep{wayEmpiricallyGroundedTechnology2022}. \end{itemize}

This appendix illustrates how partial switching from gas-based electricity to renewable electricity in the manufacturing sector affects sectoral and aggregate productivity, under three energy-accounting approaches: final, useful, and primary.  We use the primal--dual total factor productivity (TFP) framework
\begin{subequations}
\begin{align}
\hat{A}_i  &\equiv  
\underbrace{\hat{X}_i  -  \sum_{j=1}^N a_{ji} \hat{X}_{ij}  -  \sum_{f=1}^F \tilde{l}_{fi} \hat{L}_{fi}}_{\text{primal}}
\label{eq:primal_appendix}\\
 &= 
\underbrace{-\left( \hat{p}_i  -  \sum_{j=1}^N a_{ji} \hat{p}_j  -  \sum_{f=1}^F \tilde{l}_{fi} \hat{w}_f \right)}_{\text{dual}}
\label{eq:dual_appendix}
\end{align}
\end{subequations}

\noindent The variables $a_{ji}$ are the intermediate-input cost shares (so $\sum_j a_{ji} + \sum_f \tilde{l}_{fi}=1$), and $\tilde{l}_{fi}$ are the labour-cost shares.

We consider that the only inputs to the energy sector, oil, or electricity, are only capital and labour. The manufacturing sector is using labour, capital and energy (either one input or electricity or electricity and oil).

\tableofcontents


\paragraph{Why the cost shares $a_{ji}$ remain constant}

As in Appendix \ref{appendix_useful_energy}, the expenditures shares of energy remains constant when considering useful, final or primary energy as the total energy sales remain constant.

\subsection{Set-up}

We can express energy three ways: primary, final or useful. $\beta$ is the conversion factor from gas to electricity, and $\alpha$ is the conversion factor between final electricity and useful energy,  both under the same constraint: $0 \leq \alpha, \beta\leq 1$. Primary is the initial state of the energy, final is the form that is actually used. So a quantity $\beta X_G$ of final gas-generated electricity corresponds to a primary quantity of energy $X_G$ of gas. There is no difference between primary and final energy for renewables. 



We consider three industries:
\begin{itemize}
\item \textbf{G}: Gas sector, extracts raw gas and can sell electricity generated from it,
\item \textbf{R}: Renewable energy, produces renewable electricity,
\item \textbf{M}: Manufacturing, needs electricity to run its processes.
\end{itemize}


One MWh of electricity produced by renewables contains 1 MWh of final and primary energy, and provides $\alpha$ MWh of useful energy. One MWh of gas, provides $\beta$ MWh of electricity, and $\alpha\beta$ MWh of useful energy.
 

We consider three scenarios where there is a small substitution of gas by renewables: a fraction $\delta$ of M's gas consumption is replaced by renewable electricity, holding M's total output constant $\hat{X}_M=0$. Let $\gamma$ be the relative increase in renewables final electricity output.
\begin{itemize}
\item \textbf{Primary}: The manufacturing sector initially buys $X_G$ MWh of gas from G (using it directly or transforming it in electricity via direct gas combustion on-site) and buys $X_R$ MWh electricity from R (We either consider two different sectors R and G in section \ref{subappendix_primary_energy} or a single energy sector E in section \ref{subappendix_primary_energy_two_sectors}).
\item \textbf{Final}: The gas sector directly sells gas-generated electricity. M buys final energy electricity $\beta X_G$ from the gas sector G and some final renewable electricity  $X_R$  from R (We either consider two different sectors R and G in section \ref{subappendix_final_energy} or a single energy sector E in section \ref{subappendix_final_energy_two_sectors})
\item \textbf{M}: The gas sector directly sells useful gas-generated electricity. Manufacturing M buys directly useful energy gas-generated electricity $\alpha \beta X_G$ from the gas sector and some useful electricity $\alpha X_R$ from the renewables sector 
(We either consider two different sectors R and G in section \ref{subappendix_useful_energy} or a single energy sector E in section \ref{subappendix_final_energy_two_sectors}).
\end{itemize}



 
 
 Table \ref{tab_summary_energy_prod} summarises the changes in supply and demand of energy.
 Important note (\ER{VERY IMPORTANT}), what we denote here as "final energy scenario" is the scenario where all the gas output is used to produce electricity, and some is switch from gas to renewables. In the CCC scenarios we have both i) gas-electricity switched to renewables electricity (covered by "Final energy scenario" in this note), ii) gas directly used by the manufacturing switch to renewable electricity (covered by "Primary energy scenario" in this note). \er{In other words, it is equivalent to consider the electrification of} 

 

 Table \ref{tab_summary_energy_prod} summarises the changes in supply and demand of energy.

\begin{table}[H]
\centering
\begin{tabular}{@{}c|cc|cc@{}}
\toprule
        & \multicolumn{2}{c|}{Before the switch} & \multicolumn{2}{c}{After the switch} \\ \midrule
        & Gas             & Renewable            & Gas            & Renewable           \\ \midrule
Primary &     $X_G$   &    $X_R$    &  $X_G(1-\delta)$  & $X_R+\beta\delta X_G=X_R(1+\gamma)$               \\
Final   &      $\beta X_G$              &  $X_R$                      &          $\beta X_G (1-\delta)$        &        $X_R+\beta\delta X_G=X_R(1+\gamma)$               \\
Useful&  $\alpha\beta X_G$              & $\alpha X_R$                      &          $\alpha\beta X_G (1-\delta)$        &        $\alpha(X_R+\beta\delta X_G)=\alpha X_R(1+\gamma)$                   \\ \bottomrule
\end{tabular}
\caption{Production of energy}
\label{tab_summary_energy_prod}
\end{table}


The prices of the final, useful, and primary energy from the gas sector and renewables sector are noted
\[p_G^f, p_R^f,p_G^u, p_R^u,p_G^p, p_R^p. \]
The cost of final electricity is the same whether it comes from renewables or gas, 
\begin{equation}
\label{eq:price_elec}
\begin{cases}
     p_G^f&=p_R^f=p_R^p\\
      p_G^u&=p_R^u=\frac{p_R^f}{\alpha}\\
    p_G^f&=\frac{p_G^p}{\beta}
\end{cases}
\end{equation}


\subsection{Results summary}

In this paper, we focus on final energy, but within a framework similar to that of primary energy, where the energy sector delivers both electricity and fossil fuels (gas/oil). Note that in this appendix, ``final energy'' refers to final electricity. Electrification of energy productions can have diverse impacts on productivity at the sector level depending whether we are in a national account framework or in our $\zeta$ framework (\ER{Nom non contractuel}) developed for this paper, in a two or three-sector economy, or considering primary, final or useful energy.

The first result is that when considering final or useful energy, electrification does not alter the aggregate productivity in all cases (Table \ref{table_summary_results_energy}). 

The only problematic case arises when primary energy is analysed within a broad energy sector that includes both electricity and fossil fuels in our $\zeta$ framework. In this case, studying the energy sector alone correctly captures the sector's productivity growth. However, the aggregate productivity growth is non-zero and does not align with any of the others cases. 

The conclusion of Appendix (\er{A7}) where we consider productivity changes from an upstream sector does not hold. In Appendix (\er{A7}), the relationship $\hat{X}_M = \frac{\lambda_U}{\lambda_M} \hat{X}_U$ ensures that changes in the energy sector’s output are reflected in the downstream manufacturing sector. This mechanism does not hold for the energy sector itself, because useful energy remains unchanged—making it the only sector where a reduction in input can still provide exactly the same service.

\begin{table}[H]
\centering
\begin{tabular}{llll|lll}
\toprule
              & \multicolumn{3}{c}{National accounting} & \multicolumn{3}{c}{Framework} \\ \midrule
              & Primary      & Final      & Useful     
              & Primary   & Final   & Useful  \\ 
%\rowcolor[HTML]{D0CECE} 
\midrule
\textbf{\textit{Three sectors}} &              &            &             &           &         &         \\\midrule
Gas           & 0         & 0       &  0          
&   $\zeta \ln(1-\delta)$       &     $\zeta \ln(1-\delta)$      &    $\zeta \ln(1-\delta)$         \\
Renewables    & 0         & 0       &  0           
&      $\zeta \ln(\gamma)$       &       $\zeta \ln(\gamma)$     &   $\zeta \ln(\gamma)$        \\
Manufacturing & 0         & 0       &  0           
&     0     &  0        &     0    \\
Aggregate     & 0         & 0       &  0         
&    0      & 0         &   0      \\ 
%\rowcolor[HTML]{D0CECE} 
 \midrule
\textbf{\textit{Two sectors}}   &              &            &             &           &         &         \\\midrule
Energy        &    $\hat{A}_E=\hat{X}_E=-\hat{p}_E$    & 0       &  0          
&       $\hat{A}_E=\hat{X}_E=-\hat{p}_E$    &   0      &   0      \\
Manufacturing &  $\hat{A}_M=-a_{EM}\hat{X}_E$         & 0       &  0         
&      0     &    0     &       0  \\
Aggregate    & 0         & 0       &  0         
&     $\hat{A}=\frac{p_EX_E}{Y}\hat{X}_E$      &   0      &  0  \\
\bottomrule
\end{tabular}
\caption{Summary results}
\label{table_summary_results_energy}
\end{table}

Conclusion: If we adopt the final energy approach while including primary fossil fuels -— considering both the energy sector and the downstream sectors -— we encounter a quantification challenge. To ensure consistency, all energy flows would need to be converted into either final electricity or useful energy.\footnote{
\ER{François: Note that in the case displayed in the TPI version of the paper, the productivity growth we compute is the \emph{true} productivity growth of the energy sector computed with the national accounting primal/dual approach. But it would not be consistent to use it with other sectors.}
}

\textbf{The correct approach, when analysing a broad energy sector is to use the useful energy metric, ensuring that pure electrification does not affect productivity change in any sector.}

\ER{François: in other words, we will be fine with the useful energy approach within our framework.}


\subsection{Why is it different from previous appendix? (months ago)}
\ER{Subsection to be deleted after François' reading}

NB: in the previous appendix, why did I find it didn't change anything ? Because I was using our framework approach for the energy sector (with the subcase of a full switch, thus $\delta=1$) 
\[ \hat{A}_E=\hat{X}_E=\ln(\delta \beta )=\ln(\beta ),\]
However, for the manufacturing sector I was using the true productivity measure, 

  \begin{equation*}
    \left\{
    \begin{aligned}
        \hat{\Sigma}_X &= w_{EM}\hat{X}_{ME} 
                      = w_{EM}\ln \beta , \\
        \hat{\Sigma}_p &= w_{EM}\hat{p}_E
                      = w_{EM}\bigl(-\ln \beta \bigr).
    \end{aligned}
    \right.
\end{equation*}

where $w_{EM}$ is the energy share in M's cost.  Assuming $\hat{X}_M = 0$, one derives $\hat{p}_M = 0$ and
\begin{equation*}
    \hat{A}_M = -w_{EC}\ln \beta .
\end{equation*}
And then, in the aggregate, it cancels out. 
\[\hat{A}=0\]

But if we use the framework approach, then because $\hat{C}_M=0$\footnote{Because$w_{EC}$  remains constant, as the manufacturing sector's total energy bill does not change: 
\begin{equation*}
     p_GX_G 
         = \frac{p_R}{\beta }\beta X_G 
         = p_RX_R.
\end{equation*} Hence E's share in M's costs stays the same before and after.
}

\[
\left\{
\begin{aligned}
\hat{X}_M&=0\\
\hat{p}_M &= \hat{C}_M-\hat{X}_M= 0\\
\hat{A}&=0\\
\end{aligned}
\right.
\]

\newpage





\section{Three sectors: National Accounting and Our Framework}


In the following section we develop the computations of the results available in Table \ref{table_summary_results_energy}. 

We consider 3 sectors: Gas, Renewables and Manufacturing. In the next section we consider only 2: Energy and Manufacturing.

\subsection{Case 1: Primary-Energy Approach}
\label{subappendix_primary_energy}
\subsubsection{National Accounting Derivation}

In this approach, energy is measured in primary form. After a partial switch, the manufacturing sector M decreases its gas demand from $X_G$ to $(1-\delta)X_G$, while its use of renewables increases from $X_R$ to $X_R + \delta \beta X_G = X_R(1+\gamma)$.

\paragraph{Gas industry G}
Gas output shrinks by the factor $(1-\delta)$:
\[
    \hat{X}_G^p
= \ln\bigl((1-\delta) X_G\bigr)  -  \ln(X_G)
= \ln(1-\delta).
\]
From Eq.~\ref{eq:primal_appendix}, we assume all inputs to G scale down proportionally, so each input grows by $\ln(1-\delta)$. Using the same argument in the dual approach (Eq.~\ref{eq:dual_appendix}) would give $\hat{p}_G=0$, hence
\[
    \hat{A}_G
= \ln(1-\delta)
  - \sum_{j=1}^N a_{jG}\ln(1-\delta)
  - \sum_{f=1}^F \tilde{l}_{fG}\ln(1-\delta)
= \ln(1-\delta)  -  \ln(1-\delta)
= 0.
\]

\paragraph{Renewables industry R}
The renewables industry's output of final electricity rises from $X_R$ to $X_R(1+\gamma)$, so
\[
    \hat{X}_R^p
= \ln(1+\gamma).
\]
Similarly, scaling all inputs proportionally yields
\[
    \hat{A}_R = 0.
\]

\paragraph{Manufacturing sector M}
Its gas input growth is $\hat{X}_G^p=\ln(1-\delta)$, and its renewables input growth is $\hat{X}_R^p=\ln(1+\gamma)$. With $\hat{X}_M=0$ and no change in labour usage, the primal approach (Eq.~\ref{eq:primal_appendix}) gives
\[
    \hat{A}_M
= 0
  - \bigl[a_{GM}\ln(1-\delta)
     + a_{RM}\ln(1+\gamma)\bigr]
  - 0,
\]
which is equal to 0. 

We verify this by looking at M's expenditures. Since M's final energy requirement and the price of electricity remain unchanged, the total energy cost before and after the switch satisfies \(C_M = C_M'\). Let \(C_G = p_G^pX_G\) and \(C_R = p_R^pX_R\), so \(C_M = C_G + C_R\). After the switch,
\[
   C_G' = C_G(1-\delta), \quad
   C_R' = C_R(1+\gamma),\]
   Hence
\begin{equation}
\label{eq_cost_argument_M}
\begin{aligned}
         C_M'&=C_G'+C_R'=C_G  (1-\delta)+C_R(1+\gamma)\\&
     =p_G^pX_G  (1-\delta)+p_R^pX_R\left(1+\beta\delta\frac{X_G}{X_R}\right)\\
    &= p_G^pX_G  (1-\delta)+p_R^pX_R+\frac{p_G^p}{\beta}\beta \delta X_G\\
     &= p_G^pX_G +p_R^pX_R
     =C_M.
\end{aligned}
\end{equation}

It follows that
\begin{equation}
 \label{eq_weighted_share_change_cost}
\begin{aligned}
a_{GM} \ln(1-\delta) + a_{RM} \ln(1+\gamma)&= 
a_{GM}\left[ \ln(C_G') - \ln(C_G) \right] + 
a_{RM} \left[ \ln(C_R') - \ln(C_R) \right]\\&= 
\frac{C_G}{C_G + C_R} \hat{C}_G' + 
\frac{C_R}{C_G + C_R} \hat{C}_R'\\
  &=\widehat{(C_G+C_R)}\\
 &=   \ln(C_G' + C_R') - \ln(C_G + C_R) \\
&=\ln(T') - \ln(T)  \\
&= 0.
\end{aligned}
\end{equation}
because the total cost $C_M$ before and after the switch remains the same,\footnote{
Noting that the log growth of $x+y$, $\widehat{x+y}$, is $\frac{x}{x+y} \hat{x} + \frac{y}{x+y} \hat{y}$.
}
 hence
\[
    \hat{A}_M=0.
\]

\paragraph{Aggregate}
With Domar weights $\lambda_G,\lambda_R,\lambda_M$, aggregate productivity growth is
\[
    \hat{A}
= \lambda_G\hat{A}_G
  + \lambda_R\hat{A}_R
  + \lambda_M\hat{A}_M
= 0.
\]

\subsubsection{Framework Derivation}

Here, we compute industry-level productivity growth using
\[
    \hat{A}_i
= \zeta \hat{X}_i  +  (\zeta-1) \hat{p}_i.
\]

\paragraph{Gas and Renewables}
Since gas output shrinks by $\ln(1-\delta)$ while $\hat{p}_G=0$, we get
\[
    \hat{X}_G^p = \ln(1-\delta),
    \quad
    \hat{A}_G = \zeta  \hat{X}_G^p = \zeta \ln(1-\delta).
\]
Similarly, renewables output rises by $\ln(1+\gamma)$ with $\hat{p}_R=0$, so
\[
    \hat{A}_R
= \zeta \ln(1+\gamma).
\]

\paragraph{Manufacturing M}
Manufacturing output does not grow ($\hat{X}_M=0$). From the cost argument (Eq. \ref{eq_cost_argument_M}), the price $\hat{p}_M$ also stays at zero. Thus
\[
    \hat{A}_M = 0.
\]

\paragraph{Aggregate}
Finally, the Domar-weighted sum gives
\[
    \hat{A}
= \lambda_G \zeta \ln(1-\delta)
   + 
  \lambda_R \zeta \ln(1+\gamma)
= 0.
\]

\subsection{Case 2: Final-Electricity Approach}
\label{subappendix_final_energy}
\subsubsection{National Accounting Derivation}

After the partial switch, M's demand for gas-based electricity decreases by the factor $(1-\delta)$, and its demand for renewables rises by a factor $(1+\gamma)$, with exactly the same substitution condition $\beta \delta X_G = X_R \gamma$. All inputs of both G and R are scaled proportionally to their respective outputs.

\paragraph{Gas industry G}
\[
    \hat{X}_G^f
= \ln\bigl(\beta(1-\delta)X_G\bigr)
  - \ln(\beta X_G)
= \ln(1-\delta),
  \quad
  \hat{A}_G=0.
\]

\paragraph{Renewables industry R}
\[
    \hat{X}_R^f = \ln(1+\gamma),
    \quad
    \hat{A}_R = 0.
\]

\paragraph{Manufacturing M}
The price of electricity remains unchanged regardless of generation method, so substituting renewables for gas does not alter cost shares or labour usage.\footnote{For the sake of completeness, one can explicitly derive this result using the same method as in the previous case. With $\hat{X}_M=0$ and no change in M's labour usage} Hence
\[
    \hat{A}_M=0.
\]

\paragraph{Aggregate}
Again,
\begin{equation}
      \hat{A}
= \lambda_G\hat{A}_G
  + \lambda_R\hat{A}_R
  + \lambda_M\hat{A}_M
=0.
\end{equation}
\subsubsection{Framework Derivation}

Since $\hat{X}_G^f=\ln(1-\delta)$, $\hat{p}_G=0$, $\hat{X}_R^f=\ln(1+\gamma)$, $\hat{p}_R=0$, we obtain
\[
    \hat{A}_G = \zeta \ln(1-\delta),
    \quad
    \hat{A}_R = \zeta \ln(1+\gamma).
\]
Because the electricity price does not change, in manufacturing $\hat{X}_M=0$ and $\hat{p}_M=0$, so
\[
    \hat{A}_M=0,
\]
Hence, aggregate productivity growth is
\begin{equation}
    \hat{A}= \lambda_G \zeta \ln(1-\delta) + \lambda_R \zeta \ln(1+\gamma)
=0.
\end{equation}

\subsection{Case 3: Useful-Energy Approach}
\label{subappendix_useful_energy}
\subsubsection{National Accounting Derivation}

Now G and R directly sell \emph{useful} electricity to M. The partial switch replaces $\delta \alpha\beta X_G$ gas-based electricity by $\delta \alpha\beta X_G$ renewables-based electricity, leaving total useful energy demand unchanged. Both industries scale their inputs proportionally, so
\[
    \hat{X}_G^u
= \ln(1-\delta),
    \quad
    \hat{X}_R^u
= \ln(1+\gamma),
    \quad
    \hat{A}_G=0,
    \quad
    \hat{A}_R=0.
\]
In manufacturing, $\hat{X}_M=0$ and labour is unchanged, so
\[
   \hat{A}_M
= - \bigl[a_{GM}\ln(1-\delta)
       + a_{RM}\ln(1+\gamma)\bigr]
=0.
\]
Hence, 
\[
   \hat{A}=0.
\]

\subsubsection{Framework Derivation}

Exactly the same substitution argument applies; we find
\[
    \hat{X}_G^u
= \ln(1-\delta), \hat{p}_G=0,
    \quad
    \hat{X}_R^u
= \ln(1+\gamma), \hat{p}_R=0,
\]
and so
\[
    \hat{A}_G=\zeta \ln(1-\delta),
    \quad
    \hat{A}_R=\zeta \ln(1+\gamma).
\]
Since M's total energy and price do not change,
\[
    \hat{X}_M=0,\quad\hat{p}_M=0,\quad \hat{A}_M=0.
\]
Domar-aggregating again gives $\hat{A}=0$.

\section{Two sectors: National Accounting and Our Framework}


In this section we develop the computations of the results available in Table \ref{table_summary_results_energy} with only 2 sectors: Energy and Manufacturing.

\subsection{Case 4: Primary-Energy with Two Industries}
\label{subappendix_primary_energy_two_sectors}
\subsubsection{National Accounting Derivation}

Here, an ``energy industry E'' delivers two types of primary outputs to manufacturing, gas-based and renewables-based, whose total final use remains constant in terms of useful output, but not necessarily in the exact energy content. Denote $X_E^p=X_G^p + X_R^p$ in primary terms and assume $p_R^p = p_G^p/\beta$. The partial switch reduces $X_G^p$ from $X_G$ to $(1-\delta)X_G$ and increases $X_R^p$ from $X_R$ to $X_R + \beta\delta X_G$, which keeping total energy sales $C_E$ constant. 
We find
\begin{equation}
\label{eq_constant_CE}
    \begin{aligned}
    \hat{C}_E&=\ln(p_G(1-\delta)X_G+p_R(X_R+X_G\delta \beta ))-\ln(p_GX_G+p_RX_R)\\
    &=\ln(p_GX_G+p_G\frac{X_R}{\beta })-\ln(p_GX_G+p_G\frac{X_R}{\beta })\\
&=0.
\end{aligned}
\end{equation}
Hence,
\begin{equation}
\label{eq_pe_xe}
     \hat{C}_E
= \hat{X}_E^p + \hat{p}_E^p
= 0
\quad\Longrightarrow\quad
    \hat{p}_E^p
= - \hat{X}_E^p.
\end{equation}

From the primal and dual definition of productivity growth
\begin{equation}
    \hat{A}_E = - \hat{p}_E^p =\hat{X}_E^p.
\end{equation}


In manufacturing, $p_EX_E^p + \sum_f w_fL_f = p_M X_M$. Since the cost of energy is constant (Eq. \ref{eq_constant_CE}), the cost share of energy in M's expenses is constant, so
\begin{equation*}
        \hat{p}_M=\hat{C}_M-\hat{X}_M
= 0
\end{equation*}

The manufacturing sector uses energy as an intermediate input and factors. 
From the primal definition of productivity, Eq. \ref{eq:primal_appendix},\footnote{The dual definition Eq. \ref{eq:dual_appendix} gives the same result.}
\begin{equation*}
   \hat{A}_M= \hat{X}_M-a_{EM}\hat{X}_E^p= -a_{EM}\hat{X}_E^p.
\end{equation*}

Finally, the Domar-weighted sum over E and M cancels out to give
\begin{equation}
    \begin{aligned}
    \hat{A}&=\lambda_E \hat{A}_E+\lambda_M \hat{A}_M\\
    &=\frac{X_Ep_E}{Y} \hat{X}_E^p+ \frac{p_XX_M}{Y} (-a_{EM}\hat{X}_E^p)\\
    &=\frac{X_Ep_E}{Y} \hat{X}_E^p - \frac{p_XX_M}{Y} \frac{X_Ep_E}{p_XX_M}\hat{X}_E^p\\
     &=\frac{X_Ep_E}{Y} \hat{X}_E^p - \frac{X_Ep_E}{Y}\hat{X}_E^p\\
     &=0.
\end{aligned}
\end{equation}

\subsubsection{Framework Derivation}

Since, $\hat{p}_E= - \hat{X}_E^p.$ (Eq. \ref{eq_pe_xe}),
\begin{equation*}
      \hat{A}_E
= \zeta \hat{X}_E^p + (\zeta-1) \hat{p}_E
= \zeta \hat{X}_E^p + (\zeta-1)\bigl(-\hat{X}_E^p\bigr)
= \hat{X}_E^p.
\end{equation*}

In manufacturing, $\hat{X}_M=0$, $\hat{p}_M=0$, thus
\[
    \hat{A}_M=0.
\]
Hence the Domar-weighted sum is
\[
    \hat{A}
= \lambda_E \hat{A}_E
  + \lambda_M \hat{A}_M
= \lambda_E \hat{X}_E^p
\]


\subsection{Case 5: Final- or Useful-Energy Approach with Two Industries}
\label{subappendix_final_energy_two_sectors}
\subsubsection{National Accounting Derivation}

Now E delivers \emph{final} (or \emph{useful}) electricity from both gas and renewables, and M's demand for total final (or useful) energy is constant across the partial switch. The price of final energy is total expenditure on energy divided by total units of energy. Since (i) total energy use is unchanged and (ii) total expenditure on energy is also unchanged (the decline in gas spending exactly matches the rise in renewables spending), the overall price of energy must stay constant. 
Therefore,
\[
    \hat{X}_E^f
= 0,
    \quad
    \hat{p}_E= 0,
\]
hence
\[
    \hat{A}_E=0.
\]
Since the manufacturing sector faces the same final-energy demand at the same energy price,
\[
    \hat{A}_M=0,
\]
and thus $\hat{A}=0$.

\subsubsection{Framework Derivation}

If the energy industry changes its mix of inputs but the total final (or useful) output is unchanged, then
\[
   \hat{X}_E^f = 0,
   \quad
   \hat{p}_E = 0
    \Longrightarrow 
   \hat{A}_E=0.
\]
In manufacturing, neither quantity nor price changes, so $\hat{X}_M=0$ and $\hat{p}_M=0$, yielding $\hat{A}_M=0$. Domar-weighting gives $\hat{A}=0$ overall.



\subsection{Price dynamics - Two sectors, primary energy}

Now we consider that the price of gas and renewable are not fixed.

\bigskip

\subsubsection{ National Accounting Approach}

\paragraph{Energy Sector \((E)\)}

\noindent
\textbf{General formula for the log-growth of }\(p_E\):
\[
  p_E 
   =  
  \frac{p_G X_G  +  p_R X_R}{ X_E },
  \quad
  \hat{p}_E  =  \ln\!\biggl(\frac{p_E(t+1)}{ p_E(t)\!}\biggr).
\]
By taking the logarithmic derivative of the numerator minus that of the denominator, one obtains:
\[
  \hat{p}_E
   = 
  \underbrace{
    \widehat{\bigl(p_G X_G + p_R X_R\bigr)}
  }_{\text{log-growth of the sum}}
   - 
  \underbrace{
    \hat{X}_E
  }_{\text{log-growth of output}}
   .
\]
We get
\[
    \hat{p}_E
     = 
    \frac{p_G X_G}{p_E X_E} 
      (\hat{p}_G + \hat{X}_G)
     + 
    \frac{p_R X_R}{p_E X_E} 
      (\hat{p}_R + \hat{X}_R)
     - 
    \hat{X}_E.
\]

From the dual,
\[
  \hat{A}_E
   = 
  - \Bigl[
        \hat{p}_E
         \bigr)
      \Bigr],
\]
Because we only consider one sector: E. And the price of the factors don't change.

\begin{align*}
     \hat{A}_E
   &= 
  - \Bigl[
        \frac{p_G X_G}{p_E X_E} 
      (\hat{p}_G + \hat{X}_G)
     + 
    \frac{p_R X_R}{p_E X_E} 
      (\hat{p}_R + \hat{X}_R)
     - 
    \hat{X}_E
\end{align*}


\paragraph{Manufacturing Sector \((M)\)}

Manufacturing $M$ buys $X_E$ at price $p_E$, plus factors (labour, capital) whose usage and prices do not change. Let $a_{EM}$ be the cost share of energy. Then:

\paragraph{Primal.}
\[
  \hat{A}_M
   = 
  \hat{X}_M
   - 
  \bigl[a_{EM} \hat{X}_E
     + 
    \sum_{f}\tilde{l}_{fM} \hat{L}_{fM}\bigr].
\]
Since the factors do not change ($\hat{L}_{fM}=0$), this reduces to
\[
  \hat{A}_M
   = 
  \hat{X}_M 
   - 
  a_{EM} \hat{X}_E.
\]
If $M$'s output also stays unchanged ($\hat{X}_M=0$), then $\hat{A}_M=- a_{EM} \hat{X}_E$. 

From the dual,
\[
  \hat{A}_M
   = 
  - \Bigl[ 
      \hat{p}_M
       - 
      a_{EM} \hat{p}_E
    \Bigr],
\]
again ignoring factor-price terms that are zero when $\hat{w}_f=0$.

\begin{equation*}
      \hat{p}_M=\hat{C_M}=\frac{p_EX_E}{p_MX_M}\hat{C_E}=\frac{p_EX_E}{p_MX_M}\left(\hat{X}_E+\hat{p}_E\right)
\end{equation*}

\paragraph{Aggregate}

Let $\lambda_E$ and $\lambda_M$ be the Domar weights. Then
\[
  \hat{A}
   = 
  \lambda_E \hat{A}_E
   + 
  \lambda_M \hat{A}_M.
\]
Explicitly, using dual forms:
\begin{align*}
     \hat{A}
   &= 
  - \Bigl[
    \lambda_E \bigl(\hat{p}_E\bigr)
     + 
    \lambda_M \bigl(\hat{p}_M - a_{EM} \hat{p}_E\bigr)
  \Bigr]\\
 &=  - \Bigl[\frac{p_EX_E}{Y} \hat{p}_E
     + 
   \frac{p_MX_M}{Y}\bigl(\hat{p}_M - \frac{p_EX_E}{p_MX_M} \hat{p}_E\bigr)
  \Bigr]\\
  &=  - \Bigl[\frac{p_EX_E}{Y} \hat{p}_E
     + 
   \frac{p_MX_M}{Y}\bigl(\frac{p_EX_E}{p_MX_M}\left(\hat{X}_E+\hat{p}_E\right) - \frac{p_EX_E}{p_MX_M} \hat{p}_E\bigr)
  \Bigr]\\
  &= - \frac{p_EX_E}{Y}\Bigl[\hat{p}_E 
  + \left(\hat{X}_E+\hat{p}_E\right) - \hat{p}_E
  \Bigr]\\
    &= - \frac{p_EX_E}{Y}\Bigl[
  \left(\hat{X}_E+\hat{p}_E\right)
  \Bigr]\\
  &=
  - \frac{p_EX_E}{Y}\Bigl[
  \left(\hat{X}_E+\frac{p_G X_G}{p_E X_E} 
      (\hat{p}_G + \hat{X}_G)
     + 
    \frac{p_R X_R}{p_E X_E} 
      (\hat{p}_R + \hat{X}_R)
     - 
    \hat{X}_E.\right)
  \Bigr]\\
   &=
  - \frac{p_EX_E}{Y}\Bigl[
  \left(\frac{p_G X_G}{p_E X_E} 
      (\hat{p}_G + \hat{X}_G)
     + 
    \frac{p_R X_R}{p_E X_E} 
      (\hat{p}_R + \hat{X}_R)
    \right)
  \Bigr]\\
\end{align*}

Only the decrease in the price of energy appear, as if the electrification didn't take place (0 productivity gain).


\subsubsection{Framework Approach}

In this alternative definition, each sector $i$ has TFP growth
\[
  \hat{A}_i
   = 
  \zeta \hat{X}_i
   + 
  (\zeta - 1) \hat{p}_i,
\]
where $\zeta$ is a parameter (e.g.\ describing returns to scale).

\paragraph{Energy Sector \((E)\)}

\[
  \hat{A}_E
   = 
  \zeta \hat{X}_E 
   + 
  (\zeta - 1) \hat{p}_E.
\]

\paragraph{Manufacturing Sector \((M)\)}

\[
  \hat{A}_M
   = 
  \zeta \hat{X}_M
   + 
  (\zeta - 1) \hat{p}_M.
\]
With $\hat{X}_M=0$ 

\paragraph{Aggregate}

Again, Domar weighting:
\[
  \hat{A}
   = 
  \lambda_E \hat{A}_E
   + 
  \lambda_M \hat{A}_M,
\]
so



\begin{align*}
    \hat{A}
   &= 
  \lambda_E\bigl[\zeta \hat{X}_E + (\zeta -1) \hat{p}_E\bigr]
   + 
  \lambda_M\bigl[(\zeta -1) \hat{p}_M\bigr]\\
   &= 
  \lambda_E\bigl[\zeta \hat{X}_E + (\zeta -1) \hat{p}_E\bigr]
   + 
  \lambda_M\bigl[(\zeta -1) \frac{p_EX_E}{p_MX_M}\left(\hat{X}_E+\hat{p}_E\right)\bigr]\\
     &= 
  \lambda_E\bigl[\zeta \hat{X}_E + (\zeta -1) \hat{p}_E
   + 
(\zeta -1) \left(\hat{X}_E+\hat{p}_E\right)\bigr]\\
     &= 
  \lambda_E\bigl[(2\zeta-1) \hat{X}_E + 2(\zeta -1) \hat{p}_E \bigr]\\
\end{align*}

We get the same result as in the previous section if $\hat{X}_E=-\hat{p}_E $.


\section{Scenario 1: Greening electricity}

M consumes gas and electricity (both from gas and renewables). Replacing a fraction $\delta$ of gas elec by renewable elec. 

\subsection{One energy sector}

With one energy sector. 
This sector provides final energy in the form of gas and the form of electricity. 

Quantities
Before
\begin{equation*}
    \begin{cases}
        \text{Final: }Q_E^f=X_P+\beta X_G+X_R \\
         \text{Primary: }Q_E^p=X_P+ X_G+X_R \\
        \text{Useful: }Q_E^u=\alpha_P X_P+ \alpha_E\beta X_G+\alpha_E X_R \\
    \end{cases}
\end{equation*}

After
\begin{equation*}
    \begin{cases}
        \text{Final: }Q_E^f=X_P+\beta X_G(1-\delta)+X_R(1+\gamma)\\
         \text{Primary: }Q_E^p= X_P+ X_G(1-\delta)+X_R(1+\gamma\beta) \\
        \text{Useful: }Q_E^u=\alpha_P X_P+ \beta\alpha_E (1-\delta) X_G+\alpha_E X_R (1+\gamma)  \\
    \end{cases}
\end{equation*}

Sales
Before. Note that the approach does not change the sales.
\begin{equation*}
    \begin{cases}
        \text{Final: }Q_E^f=p_GX_P+\beta X_G\frac{p_G}{\beta}+X_Rp_R=p_GX_P+ X_Gp_G+X_Rp_R \\
         \text{Primary: }Q_E^p=p_GX_P+ p_GX_G+p_RX_R \\
        \text{Useful: }Q_E^u=\alpha_P X_P\frac{p_G}{\alpha_P}+ \alpha_E\beta X_G\frac{p_G}{\alpha_E\beta}+\alpha_E X_R\frac{p_R}{\alpha_E}=p_GX_P+ p_GX_G+p_RX_R \\
    \end{cases}
\end{equation*}

After
\begin{equation*}
    \begin{cases}
        \text{Final: }Q_E^f=X_Pp_G+\beta X_G(1-\delta)\frac{p_G}{\beta}+X_R(1+\gamma)p_R=X_Pp_G+ X_G(1-\delta)p_G+X_R(1+\gamma)p_R\\
         \text{Primary: }Q_E^p= p_GX_P+ X_G(1-\delta)p_G+X_R(1+\gamma\beta)p_R \\
        \text{Useful: }Q_E^u= p_GX_P+ (1-\delta)p_G X_G+ X_R (1+\gamma)p_R  \\
    \end{cases}
\end{equation*}


With a transfer of $\Delta$ from gas elec to renewable elec


Quantities
Before
\begin{equation*}
    \begin{cases}
        \text{Final: }Q_E^f=X_P+\beta X_G+X_R \\
         \text{Primary: }Q_E^p=X_P+ X_G+X_R \\
        \text{Useful: }Q_E^u=\alpha_P X_P+ \alpha_E\beta X_G+\alpha_E X_R \\
    \end{cases}
\end{equation*}

After
\begin{equation*}
    \begin{cases}
        \text{Final: }Q_E^f=X_P+\beta X_G+X_R + \underbrace{(-\Delta + \Delta )}_{=0}\\
         \text{Primary: }Q_E^p= X_P+ X_G+X_R +  \Delta(1-\frac{1}{\beta}) \\
        \text{Useful: }Q_E^u=\alpha_P X_P+ \alpha_EX_G+\alpha_E X_R + \underbrace{\alpha_E(-\Delta + \Delta )}_{=0}  \\
    \end{cases}
\end{equation*}

Hence, $\hat{Q}_E^f=\hat{Q}_E^u=0$, but $\hat{Q}_E^p\neq0$.


\paragraph{Sales}
In final and useful, the unit price of $\Delta$ is the same because it's electricity. Since quantities don't change, price do not either. But in the primary case
Before. Note that the approach does not change the sales.
\begin{equation*}
    \begin{cases}
        \text{Final: }C_E^f=p_GX_P+\beta X_G\frac{p_G}{\beta}+X_Rp_R=p_GX_P+ X_Gp_G+X_Rp_R \\
         \text{Primary: }C_E^p=p_GX_P+ p_GX_G+p_RX_R \\
        \text{Useful: }C_E^u=\alpha_P X_P\frac{p_G}{\alpha_P}+ \alpha_E\beta X_G\frac{p_G}{\alpha_E\beta}+\alpha_E X_R\frac{p_R}{\alpha_E}=p_GX_P+ p_GX_G+p_RX_R \\
    \end{cases}
\end{equation*}

After
\begin{equation*}
    \begin{cases}
        \text{Final: }C_E^f=X_Pp_G+X_Gp_G + X_Rp_R+\underbrace{(-\Delta\frac{p_G}{\beta} + \Delta p_R )}_{=0}\\
         \text{Primary: }C_E^p= p_GX_P+ X_Gp_G+X_Rp_R +  \underbrace{(-\Delta \frac{1}{\beta}p_G +\Delta p_R)}_{=0} \\
        \text{Useful: }C_E^u= p_GX_P+ p_G X_G+ X_R p_R+\underbrace{\alpha_E(-\Delta\frac{p_G}{\alpha_E\beta} + \Delta \frac{p_R}{\alpha_E} )}_{=0}  \\
    \end{cases}
\end{equation*}
Hence, $\hat{C}_E^f=\hat{C}_E^u=\hat{C}_E^p=0$




\paragraph{Average unit cost} Using 
$  \hat{C}_E= \hat{X}_E + \hat{p}_E$, i.e. $  \hat{p}_E=\hat{C}_E- \hat{X}_E$
\begin{equation*}
    \begin{cases}
        \text{Final: }\hat{p}_E^f=\hat{C}_E^f- \hat{X}_E^f=0\\
         \text{Primary: }\hat{p}_E^p= \hat{C}_E^p- \hat{X}_E^p=- \hat{X}_E^p \\
        \text{Useful: }\hat{p}_E^u=\hat{C}_E^u- \hat{X}_E^u=0  \\
    \end{cases}
\end{equation*}


\paragraph{Productivity accounting} 
Using the national accounting approach, from Eq.~\ref{eq:primal_appendix}, we assume all inputs to each of the three energies scale down proportionally to the output, respectively by $\gamma_P,\gamma_G, \gamma_R$. 
\[
    \hat{A}_E
= \hat{X}_E
  - \sum_{i\in \{P,G,R\}}\sum_{j=1}^N a_{ji}\ln(1-\delta_i)
  -  \sum_{i\in \{P,G,R\}}\sum_{f=1}^F \tilde{l}_{fi}\ln(1-\delta_i)
= \ln(1-\delta)  -  \ln(1-\delta)
= 0.
\]

Since in final and useful metrics, $\hat{X}_E=0$, and that following the previous section with three or two different sectors, the sum of inputs sums up to zero, then 
\begin{equation*}
    \begin{cases}
        \text{Final: }\hat{A}_E^f=0\\
         \text{Primary: }\hat{A}_E^p= -\hat{X}_E^p  \\
        \text{Useful: }\hat{A}_E^u=0  \\
    \end{cases}
\end{equation*}

For useful and final metrics, since the manufacturing sector faces the same final-energy demand at the same energy price, its productivity is unchanged, $\hat{A}_M=0$. 

The manufacturing sector uses energy as an intermediate input and factors. 
From the primal definition of productivity, Eq. \ref{eq:primal_appendix},\footnote{The dual definition Eq. \ref{eq:dual_appendix} gives the same result.}
\begin{equation*}
   \hat{A}_M= \hat{X}_M-a_{EM}\hat{X}_E^p= -a_{EM}\hat{X}_E^p.
\end{equation*}

Then using Domar aggregation to compute aggregate productivity, we get 
\begin{equation*}
    \begin{cases}
        \text{Final: }\hat{A}^f=0\\
         \text{Primary: }\hat{A}^p=\lambda_E \hat{A}_E+\lambda_M \hat{A}_M=\frac{X_Ep_E}{Y} \hat{X}_E^p+ \frac{p_XX_M}{Y} (-a_{EM}\hat{X}_E^p)=\frac{X_Ep_E}{Y} \hat{X}_E^p - \frac{p_XX_M}{Y} \frac{X_Ep_E}{p_XX_M}\hat{X}_E^p=0.\\
        \text{Useful: }\hat{A}^u=0  \\
    \end{cases}
\end{equation*}



\paragraph{Framework Derivation}

\begin{equation*}
    \begin{cases}
        \text{Final: }\hat{\tilde{A}}_E^f=\zeta \hat{X}_E^f+(\zeta-1) \hat{p}_E^f=0\\
         \text{Primary: }\hat{\tilde{A}}_E^p=\zeta \hat{X}_E^p+(\zeta-1) \hat{p}_E^p=\zeta \hat{X}_E^p+(1-\zeta) \hat{X}_E^p=\hat{X}_E^p\\
        \text{Useful: } \hat{\tilde{A}}_E^u=\zeta \hat{X}_E^u+(\zeta-1) \hat{p}_E^u=0  \\
    \end{cases}
\end{equation*}


In primary approach, the manufacturing sector pays the same energy bill ($\hat{C}_E^p$, since $\hat{X}_M=0$, then $\hat{p}_M=\hat{C}_M-\hat{X}_M= 0$.
Hence, $\hat{\tilde{A}}_M=0$

Then using Domar aggregation to compute aggregate productivity, we get 
\begin{equation*}
    \begin{cases}
        \text{Final: }\hat{\tilde{A}}^f=0\\
         \text{Primary: }\hat{\tilde{A}}^p=\lambda_E \hat{\tilde{A}}_E+\lambda_M \hat{\tilde{A}}_M=\lambda_E \hat{X}_E\\
        \text{Useful: }\hat{\tilde{A}}^u=0  \\
    \end{cases}
\end{equation*}

\end{document}